\section*{Abstract}
	The T0 model interprets cosmological redshift not as expansion but as energy loss of photons in the universal $\xi$-field. This document treats the field-theoretic energy loss mechanism and its connection to the Hubble constant. The quantitative derivation of $H_0 \approx 66.2\,$km/s/Mpc, unit conversion, frequency independence, and the resolution of the Hubble tension are centralized in Document~026 (Geometric Cosmology).


\section{Introduction: Rethinking the Hubble Constant}

	The conventional interpretation of Hubble's law assumes galaxies recede due to expanding space ($v = H_0 d$). This expansion paradigm requires 69\% dark energy, suffers from persistent measurement tensions, and fine-tuning problems.

	The T0 model offers a fundamentally different perspective: the universe is static, and the observed redshift arises from energy loss of photons as they propagate through the universal $\xi$-field. The Hubble constant thus transforms from a measure of space expansion into a characteristic energy loss rate.

	The time-energy duality ($\Delta E \cdot \Delta t \geq \hbar/2$) forbids a temporal beginning of the universe: a finite time interval from a singularity would require infinite energy uncertainty. The universe must have existed eternally, making space expansion unnecessary for explaining cosmic observations.

\section{The Universal $\xi$-Field}

	The cornerstone of the T0 model is the universal geometric constant:
	\begin{equation}
		\xi = \frac{4}{3} \times 10^{-4} = 1.3333\ldots \times 10^{-4}
	\end{equation}

	This dimensionless constant determines a characteristic energy scale:
	\begin{equation}
		E_\xi = \frac{1}{\xi} = 7500 \quad \text{(natural units)}
	\end{equation}

	The $\xi$-field permeates all of space and mediates interactions between photons and the vacuum. Its dynamics are described by the wave equation:
	\begin{equation}
		\square E_{\text{field}} = \left(\nabla^2 - \frac{\partial^2}{\partial t^2}\right) E_{\text{field}} = 0
	\end{equation}

\section{Energy Loss Mechanism}

	\subsection{Fundamental Energy Loss Equation}

	Photons lose energy through interaction with the $\xi$-field. The rate of energy loss follows the differential equation:
	\begin{equation}
		\frac{dE}{dx} = -\xi \cdot E
	\end{equation}

	The negative sign indicates energy loss; the linear dependence on $E$ ensures that the resulting redshift is frequency-independent (see Document~026 for the geometric justification).

	\subsection{Solution for Cosmological Distances}

	For cosmological observations ($\xi x \ll 1$), the perturbative solution yields:
	\begin{equation}
		E(x) = E_0 \left(1 - \xi \cdot x\right)
	\end{equation}

	The redshift is therefore:
	\begin{equation}
		z = \frac{E_0 - E(x)}{E(x)} \approx \xi \cdot x
	\end{equation}

	This relationship reproduces the observed linear Hubble law without space expansion.

\section{Connection to the Hubble Constant}

	Comparison with Hubble's law $z = H_0 d/c$ yields in natural units ($c = 1$):
	\begin{equation}
		\boxed{H_0^{\text{T0}} = E_H/\hbar, \quad E_H = E_0 \cdot \xi^{41/4} = 1.41 \times 10^{-33}\,\text{eV}}
	\end{equation}

	\textbf{Result:} $H_0^{\text{T0}} \approx 66.2\,$km/s/Mpc ($-1.9\%$ deviation from Planck). The exponent $41/4$ requires physical justification from $\xi$-field theory.

	For the complete numerical derivation, the step-by-step unit conversion (natural units $\rightarrow$ SI $\rightarrow$ km/s/Mpc), comparison with experimental measurements, and the resolution of the Hubble tension, see \textbf{Document~026} (Geometric Cosmology, Sections 4--6).

\section{Dimensional Analysis}

	Verification of dimensional consistency in natural units:

	\begin{center}
		\begin{tabular}{lccc}
			\toprule
			\textbf{Quantity} & \textbf{T0 Expression} & \textbf{Dimension} & \textbf{Status} \\
			\midrule
			Geometric constant & $\xi = 4/3 \times 10^{-4}$ & $[1]$ & \checkmark \\
			Energy scale & $E_\xi = 1/\xi$ & $[E]$ & \checkmark \\
			Energy loss rate & $dE/dx = -\xi \cdot E$ & $[E^2]$ & \checkmark \\
			Redshift & $z = \xi \cdot x$ & $[1]$ & \checkmark \\
			Hubble parameter & $H_0 = E_H/\hbar$ & $[T^{-1}]$ & \checkmark \\
			\bottomrule
		\end{tabular}
	\end{center}

\section{Theoretical Advantages}

	The reinterpretation of the Hubble constant as an energy loss rate solves several problems:

	\begin{enumerate}
		\item \textbf{Elimination of dark energy:} The apparent cosmic acceleration arises naturally from distance-dependent energy loss. No exotic energy form with negative pressure is required.

		\item \textbf{Fine-tuning problems solved:} The horizon problem vanishes since all regions have been in causal contact over infinite time. The flatness problem disappears for lack of initial conditions.

		\item \textbf{Parameter reduction:} Where $\Lambda$CDM requires six independent parameters, T0 derives all cosmological observables from the single geometric constant $\xi$.
	\end{enumerate}

\section{Conclusion}

	\begin{equation}
		\boxed{H_0^{\text{T0}} = E_H/\hbar \approx 66.2\,\text{km/s/Mpc}, \quad E_H = E_0 \cdot \xi^{41/4}}
	\end{equation}

	The universe is not expanding. The Hubble constant measures energy loss in the $\xi$-field, not recession. The detailed derivation via finite element simulation of the T0 vacuum, frequency independence as geometric proof, and the resolution of the Hubble tension are found in Document~026.

	\begin{thebibliography}{99}

		\bibitem{pascher_geometric_cosmology}
		J. Pascher, \textit{T0 Cosmology: Redshift as a Geometric Path Effect in a Static Universe}, T0 Document Series, Doc.~026.

		\bibitem{planck_2020}
		Planck Collaboration, \textit{Planck 2018 results. VI. Cosmological parameters}, Astron. Astrophys. 641, A6, 2020.

		\bibitem{riess_2022}
		A. G. Riess, et al., \textit{A Comprehensive Measurement of the Local Value of the Hubble Constant}, Astrophys. J. Lett. 934, L7, 2022.

	\end{thebibliography}
